\section{The agent-native cloud}\label{sec:agent-native}\label{sec:agent-native-cloud}

The roadmap states its destination without hedging, and it is worth quoting rather than
paraphrasing:
\begin{quote}
``The destination this roadmap builds toward is specific: a cloud that serves software agents
as first-class customers. An AI agent holding a wallet should be able to deploy compute on
Flux, pay per second, verify what it bought, and eventually do all of it privately --- no
account, no card, no human in the loop.'' \srcroadmap{flux-roadmap-public.md:7}
\end{quote}
That sentence names four capabilities: provisioning without a human, payment without an
account, verification of what was purchased, and eventually privacy over all of it. It is a
\vision statement --- the roadmap's own framing, not a description of the deployed system --- and
this section's job is to take it apart claim by claim rather than repeat it. Two of the four
capabilities already exist as running code. The other two are specified, not shipped, and this
paper's answer to each lives in \S14 and \S17 respectively. The distance between the sentence
above and the deployed system is smaller than a skeptical reader would guess, and the missing
piece is not an interface problem.

\paragraph{What an agent can already do.} A caller that never touches a browser can, today,
provision Flux capacity, hold a scoped and revocable credential while doing it, and pay for a
service with no account at all. Three shipped surfaces establish this.

First, an MCP (Model Context Protocol) server ships as part of the Deploy-from-Git tooling and
exposes \texttt{deploy\_app}, \texttt{update\_app} and \texttt{renew\_app} as tool calls an AI
agent can invoke directly \shipped
\srccode{deploy-with-git-frontend/\allowbreak{}server/\allowbreak{}mcp/\allowbreak{}router.js}{13-44}. These are not read-only
introspection calls: \texttt{deploy\_app}'s own description is ``Revalidate, Firebase-sign,
register, and test-install an Orbit app''
\srccode{deploy-with-git-frontend/\allowbreak{}server/\allowbreak{}mcp/\allowbreak{}createServer.js}{174-177}, and given a valid
Firebase bearer token an agent can drive a specification from draft to a signed, submitted,
on-chain registration transaction with no client-side wallet interaction and no human signing
step \srccode{deploy-with-git-frontend/\allowbreak{}server/\allowbreak{}mcp/\allowbreak{}createServer.js}{66-266}.

Second, and separately, \FluxEdge exposes a fully documented OpenAPI 3.1.1 interface branded
the ``FluxEdge API'' \shipped \srccode{pouwbackend/\allowbreak{}rest\_api/\allowbreak{}v1/\allowbreak{}openapi.yaml}{1-4}, plus a
dedicated cross-platform CLI, \texttt{fluxedge-cli}, whose machine-readable \texttt{-o json}
output mode is explicitly recommended for scripting
\srccode{pouwbackend/\allowbreak{}rest\_api/\allowbreak{}v1/\allowbreak{}openapi.yaml}{81-83,124-140}. Authority on this surface is a
bearer token, and the token's privileges are enumerated rather than all-or-nothing: a key can be
scoped to any subset of \texttt{start\_rental}, \texttt{stop\_rental}, \texttt{start\_deployment},
\texttt{stop\_deployment} and \texttt{get\_balance}
\srccode{pouwbackend/\allowbreak{}rest\_api/\allowbreak{}v1/\allowbreak{}openapi.yaml}{2571-2581}, enforced in the Go backend by a
function that maps the five rental and deployment routes to the privilege each requires and passes
every other route --- premium rental, deployment update and redeploy, failover --- for any valid key
\srccode{pouwbackend/\allowbreak{}rest\_api/\allowbreak{}ratelimit.go}{378-394}
\srccode{pouwbackend/\allowbreak{}rest\_api/\allowbreak{}api.go}{135-170}. The surface additionally rate-limits every
key and every source IP to two requests per second \srccode{pouwbackend/\allowbreak{}rest\_api/\allowbreak{}ratelimit.go}{41-42,79-88}
and automatically revokes a key or IP that accumulates five rate-limit violations inside a
five-minute window --- the key is disabled, its owner is emailed, and it is evicted from cache,
with no human review in the loop \srccode{pouwbackend/\allowbreak{}rest\_api/\allowbreak{}ratelimit.go}{143-153,158-179,203-222}.
This is a real, working instance of bounded and revocable \emph{action} authority; \S17 records
it precisely for that reason.

Third, \CumulusVPN pays for entitlement with a payment alone. A client generates a WireGuard
keypair on its own device and attaches an \texttt{OP\_RETURN} memo,
\texttt{CVPN1:\allowbreak base58(sha256(pubkey)[0{:}20])}, to an ordinary \Flux payment
\shipped \srccode{cumulusvpn/\allowbreak{}gateway/\allowbreak{}internal/\allowbreak{}entitle/\allowbreak{}entitle.go}{99-108}. Every gateway
independently scans the chain and derives entitlement from confirmed, correctly-memoed payments
\srccode{cumulusvpn/\allowbreak{}gateway/\allowbreak{}internal/\allowbreak{}entitle/\allowbreak{}entitle.go}{218-287}; there is no user database, no
login, and no server-issued token --- the public key the client already holds \emph{is} the
credential \srccode{cumulusvpn/\allowbreak{}gateway/\allowbreak{}internal/\allowbreak{}entitle/\allowbreak{}entitle.go}{1-19}. An autonomous payer
needs nothing this system does not already give it.

\paragraph{The four capabilities, scored against the evidence.} Read as a checklist against the
roadmap's own sentence, the honest score is two shipped and two not, and the two that are missing
are not symmetric with the two that ship.

\begin{itemize}
  \item \emph{Provision without a human} --- yes. The MCP tool surface and the \FluxEdge API
    both accept and act on a request with no human step in between, as above \shipped.
  \item \emph{Pay without an account} --- yes. \CumulusVPN's memo-derived entitlement ships, and
    direct per-deployment application payment is the same shape: a signed message and a
    confirmed payment, no account, no stored credential (\S16) \shipped.
  \item \emph{Verify what was bought} --- no. No tenant-facing attestation path exists today; a
    tenant cannot check that the host running its workload booted the image it was promised to
    boot. \S14 states the mechanism and the reason this is missing; it is not restated here
    \planned.
  \item \emph{Bounded spending authority} --- no. Nothing bounds the value a compromised bearer
    token can move. The action-scoping described above is real, but it scopes \emph{what} a
    token may do, never \emph{how much} it may spend; \S17 specifies the gap and the deployable
    construction that closes it \planned.
\end{itemize}

Two of four ship. The two that do not are exactly the two that decide whether an agent is merely
able to buy compute or can additionally be trusted with a budget. An agent that can provision and
pay but cannot verify or bound its own authority is a capable customer and an unsupervised
liability at the same time, and the paper does not paper over which one it currently is.

\paragraph{The interface, as a listing.} The clearest way to show what the machine-facing
interface actually looks like is to show one call, not to describe it. \FluxEdge's
\texttt{POST~/\allowbreak{}rental/\allowbreak{}start} takes a criteria object --- what kind of machine, not which machine
--- and returns a bound rental:

%% float=htbp keeps the request and its response on one page; as an inline listing
%% it broke across the page boundary mid-JSON.
\begin{lstlisting}[float=htbp, floatplacement=htbp,
  caption={\FluxEdge \texttt{POST /\allowbreak{}rental/\allowbreak{}start}: a representative agent-facing
request and response \srccode{pouwbackend/\allowbreak{}rest\_api/\allowbreak{}v1/\allowbreak{}openapi.yaml}{1330-1338,1358-1382}. The
request names criteria, not a specific machine, and \texttt{max\_price\_per\_hour} bounds this
one call. Nothing in the request, the response, or the bearer token that authorized the call
bounds what the same credential may spend across repeated calls --- the gap \S17
specifies.},label={lst:fluxedge-rental}]
// POST /rental/start
{"min_memory":16,"min_storage":500,"gpu":"RTX3080","nb_gpu":1,
 "region":"NorthAmerica","max_price_per_hour":0.5,
 "nb_computers":3,"premium":"none"}

// 200 OK
{"rentals":[{"rental_id":8782678,"price_per_hour":0.5,
  "status":"active",
  "computer":{"hash":"e5b37bd2716472","memory":16,
    "storage":500,"gpus":["RTX3080"],"nb_gpu":1,
    "region":"NorthAmerica","cluster_name":"us-east-1",
    "price_per_hour":0.5}}]}
\end{lstlisting}

\paragraph{What is missing for full autonomy.} Four gaps stand between the deployed interface and
an agent that never needs a human, and each is small enough to name precisely rather than wave at.

\begin{itemize}
  \item \emph{No self-service key issuance.} \FluxEdge's own OpenAPI surface documents
    consumption of a bearer token but not its issuance; the API's own quick-start documentation
    points a caller at a web dashboard to obtain one, and the dashboard path is the only one that ships
    \shipped \srccode{pouwbackend/\allowbreak{}rest\_api/\allowbreak{}v1/\allowbreak{}openapi.yaml}{61}. A human must mint the credential an
    agent then uses; no programmatic issuance endpoint exists anywhere on the documented surface
    \srccode{pouwbackend/\allowbreak{}rest\_api/\allowbreak{}v1/\allowbreak{}openapi.yaml}{463-2193}.
  \item \emph{No spend bound.} As above; \S17 specifies the construction.
  \item \emph{No verification path.} As above; \S14 specifies the construction.
  \item \emph{Silent underpayment on the direct-payment path.} A registration or update whose
    payment is short of the computed price is dropped with a log line and no rejection signal
    returned to the payer \srccode{flux/\allowbreak{}ZelBack/\allowbreak{}src/\allowbreak{}services/\allowbreak{}appMessaging/\allowbreak{}messageVerifier.js}{777,835}
    \srcintent{mechanism specification}. For a human this is a
    support ticket. For an unattended agent it is the one failure mode with no recovery path at
    all: the funds are already spent, the deployment does not exist, and nothing tells the caller
    either fact.
\end{itemize}

None of these four is a redesign. Self-service key issuance is an API endpoint behind the same
authorization logic that already exists; a spend bound has a deployable baseline that needs no
new consensus rule (\S17); a rejection signal is a message the network is already computing the
inputs for. The honest and more interesting claim is not that agent-native provisioning requires
new invention --- it is that four specific, narrow, already-scoped omissions are standing between
the interface as it exists today and the interface the roadmap's sentence describes.

\paragraph{Per-second GPU without an account.} The roadmap states the next step directly:
\begin{quote}
``Per-second GPU without an account. With FluxEdge folded into the platform, GPU capacity is
bought the same way: pay, run, stop. We are not fighting the GPU marketplaces on headline price;
we are the place an agent can rent a GPU without a human signing anything.''
\srcroadmap{flux-roadmap-public.md:79}
\end{quote}
This is \planned, gated on the FluxCloud/\FluxEdge merge \srcroadmap{flux-roadmap-public.md:79}.
\S15 has already established what matters here: GPU passthrough today is whole-device, granted
via Kubernetes' \texttt{nvidia.com/gpu} resource request and \texttt{NVIDIA\_VISIBLE\_DEVICES}
bus-index isolation, with no MIG, vGPU or time-slicing anywhere in the stack
\srccode{pouwbackend/\allowbreak{}rancher/\allowbreak{}deployment.go}{914-982} \srcdoc{reconciliation record}. The
consequence for this section, stated once and not re-derived: a per-second \emph{price} does not
require a per-second \emph{allocation} primitive, and nothing here blocks quoting or billing a
whole GPU by the second. What per-second billing cannot yet mean, on the allocation primitive
\S15 describes, is fractional-GPU multi-tenancy --- a device is sold whole or not at all
\srcinferred.

\paragraph{The composite picture.} An agent today can obtain \Flux capacity and pay for it with
no account and no human step; \FluxEdge's scoped bearer tokens and \CumulusVPN's memo-derived
entitlement are both running code, not proposals. What that same agent cannot do is confirm that
the host it paid actually booted the image it was promised, and it cannot hold a credential whose
worst-case loss is anything less than the full balance behind it. Both gaps have specifications
elsewhere in this paper --- \S14 for the first, \S17 for the second --- and \S\ref{sec:feasibility}
classes neither as Extension: attested execution is Research$\to$Engineering (row 4), and bounded
spending is Research (row 13) because the L1 cannot enforce $\budget$
(\Cref{prop:budget-balance}, \Cref{tab:two-authorities}). What is deployable today is the funded
allowance of \Cref{con:allowance} --- a funded key with a swept revocation address, which the
mechanism specification classes as Extension \srcintent{mechanism specification} --- whose
bound is the balance behind the key, not a budget the chain enforces.

\paragraph{Why this matters architecturally rather than commercially.} The properties an
autonomous customer needs --- no account, no card, no human in the loop, verifiable delivery, a
bounded worst-case loss --- are exactly the properties a careful, distrustful human customer
wants as well. A human who reads a memo-derived entitlement scheme correctly does not want a
password database holding their identity either; a human who deploys a container wants to know
what booted it just as much as an agent does; a human handing a long-lived API key to a CI
pipeline wants that key's blast radius bounded for the same reason a delegated agent needs one.
Building the interface this section describes for agents is not a separate product from building
it for people who do not want to trust the platform by default --- it is the same interface,
and this paper makes no claim about how many agents will use it.

\subsection{What such an agent would run}
\label{sec:agent-native-workload}

The interface question of this section --- how an agent obtains capacity --- is separable from what it
then does with it, but the two meet at one point worth stating. The workloads an agent runs on behalf of
an ordinary person are, in the main, small task-specialised models rather than frontier ones, and
\S\ref{sec:ac-workload-fit} argues that this class of work has no interconnect requirement and therefore
no structural reason to centralise. That is the reason the agent-native argument and the accelerated
compute argument are the same argument seen from two ends: the network that can be bought from without
an account is also the network whose dispersion stops being a handicap for the work being bought.

Two of this section's four capabilities are what make that composition real rather than rhetorical.
Payment without an account means a caller with no relationship to maintain can buy inference for one
task and stop. Verification of what was bought --- absent today (\S\ref{sec:attested-execution}) ---
matters disproportionately for a small specialised model, because a caller usually cannot tell from an
answer whether the model that produced it was the one they paid for. Neither property is peculiar to
autonomous callers; both are simply more visible when there is no human to absorb the ambiguity.

%% Drain this section's float queue before the next begins.
\FloatBarrier
